\chapter{Całki po trajektoriach i funkcjonały miary Wienera}
\label{ch:path-integrals}
\paragraph{Odsyłacze międzytomowe.} Wykorzystanie metod próbkowania do aproksymacji całek omawia tom IV, rozdział~\ref{IV-ch:monte-carlo}.


\section{Motywacja: całkujemy po przestrzeni funkcji}

W zwykłej całce Lebesgue'a argumentem funkcji podcałkowej jest liczba
albo punkt przestrzeni euklidesowej. W całce \emph{po trajektoriach}
argumentem jest cała funkcja $w:[0,T]\to\R^d$. Przykładowo
$F(w)=\exp(-\int_0^T V(w(s))\,\dd s)$ zależy od tego, w jakich
obszarach przestrzeni przebywała trajektoria w całym przedziale czasu.
Chcemy nadać ścisły sens wyrażeniu $\int F(w)\,\dd\mu(w)$, a nie
wprowadzać nieokreślony symbol $\mathcal D w$.

W rozdziale~\ref{ch:cameron-martin} skonstruowaliśmy miarę Wienera.
Teraz użyjemy jej do obliczeń, zbudujemy \emph{miarę mostu Browna},
a następnie przejdziemy od funkcjonałów trajektorii do jądra
operatora Feynmana--Kaca. Jest to punkt wyjścia do matematycznej
interpretacji niektórych całek stosowanych w mechanice kwantowej.
Przypominamy rozróżnienie: $\int_0^T H_s\,\dd W_s$ jest zmienną
losową, natomiast $\int F(w)\,\dd\mu_W(w)$ jest liczbą, o ile
funkcjonał $F$ jest całkowalny.

\begin{definicja}[Kanoniczna przestrzeń trajektorii i miara Wienera]
\label{def:path-wiener}
Niech $d\in\N$, $T>0$ i $x\in\R^d$. Oznaczmy
$\calC_{x,T}=\{w\in C([0,T];\R^d):w(0)=x\}$, z normą
supremum i borelowskim $\sigma$-ciałem. Niech $B$ będzie
$d$-wymiarowym standardowym ruchem Browna. Miara
$\mu_{x,T}:=\Law((x+B_s)_{0\le s\le T})$ na $\calC_{x,T}$
nazywa się \emph{miarą Wienera rozpoczynającą się w $x$}.
Jeżeli $F:\calC_{x,T}\to\R$ jest funkcjonałem borelowskim
i $\int|F|\,\dd\mu_{x,T}<\infty$, to jego \emph{całką
po trajektoriach} jest $\int_{\calC_{x,T}}F(w)\,\dd\mu_{x,T}(w)$.
\end{definicja}

Definicja nie używa fikcyjnego „jednostajnego rozkładu” na
nieskończeniewymiarowej przestrzeni funkcji. Miara Wienera jest
skoncentrowana na trajektoriach ciągłych, ale prawie na pewno
nigdzie nieróżniczkowalnych. Tych własności nie można ignorować
w całkach po trajektoriach.

Dla $t>0$ oznaczmy przez
$p_t(x,y)=(2\pi t)^{-d/2}\exp(-\|y-x\|^2/(2t))$
jądro cieplne generatora $\frac12\Delta$.

\begin{twierdzenie}[Całki cylindryczne i rozkłady skończeniewymiarowe]
\label{th:path-cylinder}
Niech $0<t_1<\cdots<t_n\le T$, niech
$\Phi:(\R^d)^n\to\R$ będzie borelowska i całkowalna
względem rozkładu $(x+B_{t_1},\ldots,x+B_{t_n})$.
Wówczas
\[
\begin{split}
 &\int_{\calC_{x,T}}\Phi(w(t_1),\ldots,w(t_n))\,
       \dd\mu_{x,T}(w)\\
 &\quad=\int_{(\R^d)^n}\Phi(x_1,\ldots,x_n)
   \prod_{k=1}^n p_{t_k-t_{k-1}}(x_{k-1},x_k)
   \,\dd\lambda_{dn}(x_1,\ldots,x_n),
\end{split}
\]
przy konwencji $t_0=0$, $x_0=x$. Funkcjonały cylindryczne
wyznaczają jednoznacznie całki wszystkich ograniczonych
borelowskich funkcjonałów na $\calC_{x,T}$.
\end{twierdzenie}
\begin{proof}
Przyrosty $B_{t_k}-B_{t_{k-1}}$ są niezależne i mają gęstości
$p_{t_k-t_{k-1}}(0,\cdot)$. Odwzorowanie liniowe wysyłające
przyrosty na kolejne położenia ma wyznacznik równy $1$,
więc wspólna gęstość położeń jest iloczynem zapisanym w tezie.
Otrzymujemy wzór najpierw dla funkcji wskaźnikowych
prostokątów, następnie przez liniowość i twierdzenie o
zbieżności monotonicznej dla wszystkich nieujemnych
funkcji borelowskich, a przez rozdzielenie części dodatniej
i ujemnej dla całkowalnych.

Wreszcie borelowskie $\sigma$-ciało na przestrzeni
ciągłych trajektorii jest generowane przez odwzorowania
$w\mapsto w(q)$ dla wymiernych $q\in[0,T]$.
Istotnie, norma supremum dwóch trajektorii jest supremum
odległości w chwilach wymiernych. Klasa zbiorów, dla
których miary się zgadzają, jest klasą Dynkina zawierającą
cylindry i ich skończone przecięcia; twierdzenie o
$\pi$-$\lambda$-układach daje jednoznaczność.
\end{proof}

\section{Most Browna: prawidłowe warunkowanie względem końca trajektorii}

Zdarzenie $\{x+B_T=y\}$ ma prawdopodobieństwo zero dla
każdego ustalonego $y$. Dlatego nie możemy zdefiniować mostu
przez elementarny iloraz prawdopodobieństw.
Rozwiązaniem jest \emph{regularny rozkład warunkowy}, który
możemy w tym przypadku skonstruować wprost.

\begin{definicja}[Most Browna pomiędzy ustalonymi końcami]
\label{def:brownian-bridge-path}
Niech $x,y\in\R^d$, $T>0$, a $B$ będzie standardowym
$d$-wymiarowym ruchem Browna. Proces
\[
 \beta_s^{x,y,T}=x+\frac{s}{T}(y-x)+B_s-\frac{s}{T}B_T,
 \qquad 0\le s\le T,
\]
nazywamy \emph{mostem Browna z $x$ do $y$ w czasie $T$}.
Jego prawo na
$\calC_{x,y,T}:=\{w\in C([0,T];\R^d):w(0)=x,\ w(T)=y\}$
oznaczamy przez $\mu_{x,y,T}^{\mathrm{br}}$.
\end{definicja}

\begin{twierdzenie}[Konstrukcja gaussowska i warunkowanie]
\label{th:bridge-conditional}
Dla każdego $x,y\in\R^d$ proces $\beta^{x,y,T}$ ma
ciągłe trajektorie, średnią
$m_s=x+\frac{s}{T}(y-x)$ i macierz kowariancji
\[
\operatorname{Cov}(\beta_s^{x,y,T},\beta_r^{x,y,T})
 =\left(\min(s,r)-\frac{sr}{T}\right)I_d.
\]
Dla $0<t_1<\cdots<t_n<T$ jego rozkład wektora
$(\beta_{t_1},\ldots,\beta_{t_n})$ ma gęstość
\[
\frac{p_{t_1}(x,x_1)\,p_{t_2-t_1}(x_1,x_2)
 \cdots p_{t_n-t_{n-1}}(x_{n-1},x_n)
 p_{T-t_n}(x_n,y)}{p_T(x,y)}.
\tag{15.1}\label{eq:bridge-density}
\]
Ponadto rodzina $y\mapsto\mu_{x,y,T}^{\mathrm{br}}$
jest wersją regularnego rozkładu warunkowego miary Wienera
$\mu_{x,T}$ względem $w(T)=y$: dla nieujemnego borelowskiego
funkcjonału $F$ i nieujemnej borelowskiej funkcji $g$,
\[
 \int F(w)g(w(T))\,\dd\mu_{x,T}(w)
 =\int_{\R^d}p_T(x,y)g(y)
    \left(\int F(w)\,\dd\mu_{x,y,T}^{\mathrm{br}}(w)\right)
    \dd\lambda_d(y).
\tag{15.2}\label{eq:bridge-disintegration}
\]
\end{twierdzenie}
\begin{proof}
Wektor $(B_{t_1},\ldots,B_{t_n},B_T)$ jest gaussowski.
Dla $s,r\le T$ obliczamy
$\operatorname{Cov}(B_s-\frac sT B_T,
B_r-\frac rT B_T)=\min(s,r)-sr/T$.
Co więcej, proces $B_s-sB_T/T$ jest niezależny od $B_T$:
ma zerową kowariancję z $B_T$, a cała rodzina jest wspólnie
gaussowska. Zatem warunkowanie procesu $x+B$ przy wartości
końcowej $y$ daje proces o tej samej kowariancji, lecz z
podstawioną średnią $m_s$. To jest dokładnie prawo
$\beta^{x,y,T}$.

Gęstość wspólnego rozkładu
$(x+B_{t_1},\ldots,x+B_{t_n},x+B_T)$ jest iloczynem
jąder cieplnych. Dzieląc przez dodatnią gęstość brzegową
$p_T(x,y)$, otrzymujemy \eqref{eq:bridge-density} dla
każdego $y$ (przy takim wyborze wersji jądra).
Tożsamość \eqref{eq:bridge-disintegration} wynika więc
najpierw dla cylindrycznych $F$ i prostych $g$ z twierdzenia
Fubiniego. Klasa tych funkcjonałów generuje całe borelowskie
$\sigma$-ciało przestrzeni trajektorii; twierdzenie o klasie
monotonicznej rozszerza tożsamość na wszystkie nieujemne
borelowskie $F$ i $g$. Mierzalność względem $y$ wynika
również z jawnej reprezentacji
$\beta^{x,y,T}=\beta^{x,0,T}+(s/T)y$.
\end{proof}

\begin{wniosek}[Gęstość początkowego fragmentu mostu]
\label{cor:bridge-initial-density}
Niech $0<s<T$, $y\in\R^d$ i niech $F$ będzie
ograniczonym borelowskim funkcjonałem zależnym wyłącznie
od trajektorii do chwili $s$. Wówczas
\[
 \E_{x,y,T}^{\mathrm{br}}F
 =\frac{1}{p_T(x,y)}\E_x\left[
    F(x+B)\,p_{T-s}(x+B_s,y)\right].
\tag{15.3}\label{eq:bridge-likelihood}
\]
Dla każdego $s<T$ prawo mostu na $[0,s]$ jest więc
równoważne prawu zwykłego ruchu Browna na $[0,s]$.
Na pełnym przedziale $[0,T]$ miary te są natomiast
wzajemnie singularne.
\end{wniosek}
\begin{proof}
Po ustaleniu przeszłości do $s$ niezależny przyrost
$B_T-B_s$ ma gęstość $p_{T-s}(x+B_s,\cdot)$.
Stosujemy tę własność do wspólnej miary
$\E_x[F\ind{\{x+B_T\in A\}}]$, a następnie porównujemy
gęstości względem $\lambda_d(y)$ i używamy ciągłej wersji
jądra z \eqref{eq:bridge-disintegration}. Współczynnik
$p_{T-s}/p_T$ jest ściśle dodatni i skończony,
co daje równoważność miar na $[0,s]$.
W chwili końcowej most spełnia $w(T)=y$ z
prawdopodobieństwem $1$, a zwykła miara Wienera przypisuje
temu zdarzeniu prawdopodobieństwo $0$. To dowodzi
singularności na $[0,T]$.
\end{proof}

\begin{przyklad}[Położenie w połowie drogi]
W jednowymiarowym moście z $x$ do $y$ w czasie $T$
zmienna $\beta_{T/2}$ ma rozkład
$\Normal((x+y)/2,T/4)$, a nie rozkład jednostajny.
Średnia jest liniową interpolacją końców, natomiast
wariancja $T/4$ jest konsekwencją wymagania zakończenia
trajektorii dokładnie w $y$.
\end{przyklad}



\begin{zadanie}[Prawdopodobieństwo przejścia mostu w przód]
Niech $0\le s<t<T$ i $\beta^{x,y,T}$ będzie
$d$-wymiarowym mostem Browna. Wyznaczyć gęstość warunkową
$\beta_t$ przy zadanej wartości $\beta_s=z$.

\emph{Rozwiązanie.} Z formuły
\eqref{eq:bridge-density}, po wycałkowaniu pozostałych
współrzędnych, otrzymujemy jądro
\[
 p^{\mathrm{br}}_{s,t}(z,z')=
 \frac{p_{t-s}(z,z')p_{T-t}(z',y)}{p_{T-s}(z,y)}.
\]
Mianownik jest dodatni, a całka prawej strony względem
$z'$ wynosi $1$ dzięki półgrupowej własności jądra
cieplnego. Jądro zależy od obu chwil $s,t$, zatem most
jest \emph{niejednorodnym w czasie} procesem Markowa.
Dodatkowy czynnik $p_{T-t}(z',y)$ nadaje większą wagę
położeniom, z których łatwiej dotrzeć do końca $y$.
To wyjaśnia kierunek dryfu w następnym twierdzeniu.
\end{zadanie}

\section{Dynamika mostu: dryf, który wymusza punkt końcowy}

\begin{twierdzenie}[Stochastyczne równanie mostu Browna]
\label{th:bridge-sde}
Niech $x,y\in\R$, $T>0$, a $W$ będzie standardowym
ruchem Browna. Dla $0\le s<T$ określmy
\[
 Y_s=x+\frac{s}{T}(y-x)+(T-s)\int_0^s\frac1{T-r}\,\dd W_r.
\tag{15.4}\label{eq:bridge-adapted}
\]
Proces $Y$ ma to samo prawo co most Browna z $x$ do $y$
i przyjmuje ciągłą modyfikację na $[0,T]$ spełniającą
$Y_T=y$. Na $[0,T)$ jest jedynym silnym rozwiązaniem
\[
 \dd Y_s=\frac{y-Y_s}{T-s}\,\dd s+\dd W_s,
 \qquad Y_0=x.
\tag{15.5}\label{eq:bridge-sde}
\]
\end{twierdzenie}
\begin{proof}
Na każdym $[0,T-\varepsilon]$ funkcja
$(T-s)^{-1}$ jest ograniczona i kwadratowo całkowalna,
więc całka Itô w \eqref{eq:bridge-adapted} jest
poprawnie określona i adaptowana. Z izometrii Wienera,
dla $s\le r<T$,
\[
\begin{split}
&\Cov\left((T-s)\int_0^s\frac{\dd W_u}{T-u},
 (T-r)\int_0^r\frac{\dd W_u}{T-u}\right)\\
&\qquad=(T-s)(T-r)\int_0^s\frac{\dd u}{(T-u)^2}
 =s-\frac{sr}{T}.
\end{split}
\]
Średnia $Y_s$ jest równa $x+s(y-x)/T$.
Wspólna gaussowskość daje zatem te same rozkłady
skończeniewymiarowe, co w
Twierdzeniu~\ref{th:bridge-conditional}.
Na każdym $[0,T-\varepsilon]$ proces $Y$ ma trajektorie
ciągłe. Uzasadnijmy osobno zachowanie w chwili $T$.
Wybierzmy zbiór $D=\mathbb Q\cap[0,T)$, gęsty w $[0,T)$.
Zgodność rozkładów skończeniewymiarowych $Y$ i mostu
$\beta^{x,y,T}$ oznacza zgodność ich praw na
przeliczalnej przestrzeni współrzędnych $\mathbb R^D$.
Dla mostu z ciągłości na $[0,T]$ wynika, że dla każdego
$k\in\N$ istnieje $\delta>0$ takie, że
$|\beta_q-y|<1/k$ dla wszystkich $q\in D$,
$T-\delta<q<T$. Zdarzenie to jest borelowskie
w przestrzeni $\mathbb R^D$, zatem ten sam warunek
spełnia $Y$ prawie na pewno. Ponieważ $Y$ jest ciągły
na $[0,T)$, oszacowanie rozciąga się z wymiernych $q$
na wszystkie $s$ dostatecznie bliskie $T$. Otrzymujemy
$\lim_{s\uparrow T}Y_s=y$ prawie na pewno. Ustalając
$Y_T=y$, dostajemy ciągłe przedłużenie i prawo mostu
na całym $[0,T]$. Nie wystarczyłaby tu sama zbieżność
$Y_s\to y$ według prawdopodobieństwa.

Z reguły iloczynu, ponieważ $s\mapsto T-s$ ma
skończoną wariację, otrzymujemy
\[
\dd\left((T-s)\int_0^s\frac{\dd W_r}{T-r}\right)
 =-\left(\int_0^s\frac{\dd W_r}{T-r}\right)\dd s+\dd W_s.
\]
Po dodaniu pochodnej deterministycznej średniej
upraszcza się to do \eqref{eq:bridge-sde}.
Jednoznaczność na $[0,T-\varepsilon]$ wynika z
liniowości równania lub z mnożenia różnicy dwóch
rozwiązań przez $(T-s)^{-1}$; ponieważ
$\varepsilon>0$ jest dowolne, otrzymujemy
jednoznaczność na $[0,T)$.
\end{proof}

\begin{uwaga}[Intuicja dryfu mostu]
Współczynnik $(y-Y_s)/(T-s)$ jest tym większy, im
mniej czasu pozostało do końca i im dalej proces
znajduje się od $y$. Nie jest on ograniczony przy
$s\uparrow T$, ale równanie jest dobrze określone
na każdym przedziale $[0,T-\varepsilon]$, a w granicy
trajektoria dochodzi do wyznaczonego punktu.
To bardzo konkretny przykład \emph{niejednorodnego
w czasie procesu Markowa}.
\end{uwaga}

\section{Funkcjonały liniowe i pierwsze obliczenia całek po trajektoriach}

\begin{twierdzenie}[Całka trajektorii mostu jest gaussowska]
\label{th:bridge-integrated}
Jeżeli $d=1$, $T>0$ i $x,y\in\R$, to
\[
 A_T:=\int_0^T\beta_s^{x,y,T}\,\dd s
 \sim\Normal\!\left(\frac{T(x+y)}2,\frac{T^3}{12}\right).
\]
W szczególności dla każdego $u\in\R$,
\[
 \int_{\calC_{x,y,T}}\exp\left(iu\int_0^Tw(s)\,\dd s\right)
     \,\dd\mu_{x,y,T}^{\mathrm{br}}(w)
 =\exp\left(\frac{iuT(x+y)}2-\frac{u^2T^3}{24}\right).
 \tag{15.6}\label{eq:bridge-area-char}
\]
\end{twierdzenie}
\begin{proof}
Sumy Riemanna są liniowymi kombinacjami wspólnie
gaussowskich zmiennych, a z ciągłości trajektorii zbiegają
prawie na pewno. Aby uzasadnić również zbieżność w $L^2$,
niech $\beta^{(n)}_s$ oznacza funkcję schodkową równą
$\beta_{t_k}$ na przedziale $[t_k,t_{k+1})$.
Z nierówności Cauchy'ego--Schwarza
\[
\E\left|\int_0^T(\beta^{(n)}_s-\beta_s)\,\dd s\right|^2
\le T\int_0^T\E|\beta^{(n)}_s-\beta_s|^2\,\dd s.
\]
Dla mostu
$\E|\beta_t-\beta_s|^2\le |t-s|+
\|y-x\|^2|t-s|^2/T^2$,
bo wariancja przyrostu wynosi
$|t-s|-(t-s)^2/T$; oba składniki znikają wraz
ze średnicą podziału. Prawa strona
zmierza więc do zera wraz ze średnicą podziału. Funkcje
charakterystyczne gaussowskich sum przechodzą do granicy,
co pokazuje, że sama całka ma rozkład gaussowski.
Jej średnia jest całką funkcji $m_s$ i wynosi
$T(x+y)/2$. Z twierdzenia Fubiniego
\[
 \Var(A_T)=\int_0^T\!\int_0^T
 \left(\min(s,r)-\frac{sr}{T}\right)\dd s\,\dd r
 =\frac{T^3}{3}-\frac{T^3}{4}=\frac{T^3}{12}.
\]
Wzór \eqref{eq:bridge-area-char} jest teraz funkcją
charakterystyczną rozkładu normalnego.
\end{proof}

\begin{zadanie}[Potencjał liniowy: jawna całka funkcyjna]
Dla $\lambda\in\R$ obliczyć
$\E_{x,y,T}^{\mathrm{br}}
\exp(-\lambda\int_0^T\beta_s\,\dd s)$.

\emph{Rozwiązanie.} Z twierdzenia powyżej i wzoru na
transformatę Laplace'a rozkładu normalnego otrzymujemy
\[
 \E_{x,y,T}^{\mathrm{br}}
  \exp\!\left(-\lambda\int_0^T\beta_s\,\dd s\right)
 =\exp\left(-\frac{\lambda T(x+y)}2
           +\frac{\lambda^2T^3}{24}\right).
\]
Funkcjonał jest całkowalny pomimo tego, że potencjał
$V(z)=\lambda z$ nie jest ograniczony. Wynik pochodzi
z rachunku gaussowskiego, a nie z ogólnego twierdzenia
o całkach z dowolnym nieograniczonym potencjałem.
\end{zadanie}


\section{Nieliniowy funkcjonał: całka z kwadratu ruchu Browna}

\begin{twierdzenie}[Transformata Laplace'a energii trajektorii]
\label{th:brownian-square-laplace}
Niech $B$ będzie standardowym jednowymiarowym
ruchem Browna z $B_0=0$. Dla $\alpha\ge0$ i $T>0$,
\[
 \E\exp\left(-\alpha\int_0^T B_s^2\,\dd s\right)
 =\frac1{\sqrt{\cosh(T\sqrt{2\alpha})}}.
\tag{15.7}\label{eq:brownian-square-laplace}
\]
\end{twierdzenie}
\begin{proof}
Dla $x\in\R$ rozważmy $X_s=x+B_s$ i funkcję
$u(t,x)=\E_x\exp(-\alpha\int_0^tX_s^2\,\dd s)$.
Zamiast zakładać z góry, że $u$ jest rozwiązaniem PDE,
znajdziemy jego kandydata i zweryfikujemy go przez
wzór Itô. Szukamy funkcji $v(t,x)=A(t)e^{-q(t)x^2}$
spełniającej
$\partial_tv=\tfrac12\partial_{xx}v-\alpha x^2v$,
$v(0,x)=1$. Po podstawieniu otrzymujemy
\[
 q'(t)=\alpha-2q(t)^2,\qquad
 A'(t)=-q(t)A(t),\qquad q(0)=0,\quad A(0)=1.
\]
Dla $\alpha>0$ rozwiązaniami są
\[
 q(t)=\sqrt{\alpha/2}\tanh(t\sqrt{2\alpha}),\qquad
 A(t)=\cosh(t\sqrt{2\alpha})^{-1/2}.
\]
Zatem $0<v(t,x)\le1$ oraz $v(0,x)=1$.
Dla ustalonego $t>0$ stosujemy wzór Itô do procesu
\[
 M_s=\exp\left(-\alpha\int_0^sX_r^2\,\dd r\right)
 v(t-s,X_s),\qquad 0\le s\le t.
\]
Drift znika na mocy równania PDE. Po zatrzymaniu w
$\tau_R=\inf\{s:|X_s|\ge R\}$ część stochastyczna
jest martyngałem. Ponieważ $0\le M_s\le1$, można
usunąć zatrzymanie przez zbieżność zdominowaną,
a w chwili $t$ mamy $v(0,X_t)=1$. Dostajemy
$v(t,x)=\E_xe^{-\alpha\int_0^tX_s^2\,\dd s}$.
Podstawiając $x=0$, otrzymujemy tezę. Dla
$\alpha=0$ obie strony są równe $1$.
\end{proof}

\begin{uwaga}[Dwa rodzaje funkcjonałów]
Całka $\int_0^T B_s\,\dd s$ jest gaussowska,
ponieważ jest liniowym funkcjonałem procesu
gaussowskiego. Całka $\int_0^T B_s^2\,\dd s$
nie jest gaussowska: jest nieujemna i ma transformację
Laplace'a \eqref{eq:brownian-square-laplace}.
Obliczenie drugiej wymaga więcej niż znajomości
funkcji kowariancji i prowadzi naturalnie do
równania parabolicznego.
\end{uwaga}

\section{Feynmana--Kaca półgrupa ważona i jej jądro}

\begin{definicja}[Półgrupa z potencjałem]
\label{def:path-weighted-semigroup}
Niech $V:\R^d\to\R$ będzie ograniczoną funkcją borelowską.
Dla $t\ge0$ i ograniczonej borelowskiej funkcji
$f:\R^d\to\R$ określamy
\[
 (P_t^Vf)(x)=\E_x\left[
 \exp\left(-\int_0^tV(X_s)\,\dd s\right)f(X_t)\right],
 \quad X_s=x+B_s.
\tag{15.8}\label{eq:path-FK-semigroup}
\]
Całka względem $s$ jest zwykłą całką Lebesgue'a;
całość jest całką po prawie trajektorii $X$.
\end{definicja}

\begin{twierdzenie}[Własność półgrupowa i jądro mostowe]
\label{th:path-FK-kernel}
Przy założeniach Definicji~\ref{def:path-weighted-semigroup}:
\begin{enumerate}
\item $P_{s+t}^V=P_s^VP_t^V$ dla $s,t\ge0$,
$P_0^V=I$ oraz
$\|P_t^V f\|_\infty\le e^{t\|V\|_\infty}\|f\|_\infty$;
\item $P_t^V$ jest operatorem dodatnim;
\item dla $t>0$ funkcja
\[
 k_t^V(x,y)=p_t(x,y)\,
 \E_{x,y,t}^{\mathrm{br}}\left[
   \exp\left(-\int_0^tV(\beta_s)\,\dd s\right)\right]
 \tag{15.9}\label{eq:path-FK-kernel}
\]
jest mierzalnym jądrem całkowym, tj.
$(P_t^Vf)(x)=\int_{\R^d}k_t^V(x,y)f(y)\,\dd\lambda_d(y)$;
\item jądra spełniają tożsamość Chapmana--Kołmogorowa:
\[
 k_{s+t}^V(x,y)=\int_{\R^d}k_s^V(x,z)k_t^V(z,y)
 \,\dd\lambda_d(z)
\]
dla $s,t>0$ i prawie wszystkich $y$
(przy ustalonych $x,s,t$).
\end{enumerate}
\end{twierdzenie}
\begin{proof}
Niech $X=x+B$. Dla ustalonych $s,t\ge0$ rozdzielamy
całkę potencjału na $[0,s]$ i $[s,s+t]$. Dzięki własności
Markowa, dla funkcjonału dalszej trajektorii otrzymujemy
\[
 \E_x\left[e^{-\int_s^{s+t}V(X_r)\,\dd r}
    f(X_{s+t})\mid\calF_s\right]=(P_t^V f)(X_s).
\]
Tożsamość najpierw zachodzi dla prostych funkcjonałów
zależnych od skończenie wielu chwil, a przez aproksymację
monotoniczną dla funkcjonałów mierzalnych całej przyszłej
trajektorii. Funkcjonał $\int_s^{s+t}V(X_r)\,\dd r$
jest mierzalny: można go otrzymać jako granicę całek
funkcji prostych względem czasu, ponieważ
$(r,\omega)\mapsto V(X_r(\omega))$ jest produktowo
mierzalne. Własność wieży daje $P_{s+t}^V=P_s^VP_t^V$.
Dodatniość i oszacowanie normy wynikają z granic
$e^{-t\|V\|_\infty}\le e^{-\int_0^tV}\le e^{t\|V\|_\infty}$.

Aby otrzymać jądro, stosujemy rozkład warunkowy
\eqref{eq:bridge-disintegration} do funkcjonału
$F(w)=\exp(-\int_0^tV(w(r))\,\dd r)$
i funkcji $g(w(t))=f(w(t))$ (najpierw dla funkcji
nieujemnych, potem przez rozkład na części dodatnią
i ujemną). Dostajemy \eqref{eq:path-FK-kernel}.
Złożenie dwóch operatorów całkowych, twierdzenie
Fubiniego i własność półgrupowa dają ostatnią tezę
jako równość jąder dla prawie wszystkich $y$.
\end{proof}

\begin{uwaga}[Brak konieczności istnienia gęstości trajektorii]
Wyrażenie \eqref{eq:path-FK-kernel} nie sugeruje, że
istnieje gęstość miary Wienera względem jakiejś
„miary Lebesgue'a” na $C([0,t])$. Czynnik $p_t(x,y)$
pochodzi wyłącznie z gęstości \emph{końcowego położenia}.
Pozostały czynnik jest zwykłą wartością oczekiwaną po
moście, który ma dobrze określoną miarę probabilistyczną.
\end{uwaga}

\section{Dyskretyzacja czasu: dlaczego powstaje całka po trajektoriach?}

\begin{twierdzenie}[Skończeniewymiarowe aproksymacje całki po trajektoriach]
\label{th:path-time-slice}
Niech $T>0$, $x\in\R^d$, $V:\R^d\to\R$ będzie
\emph{ograniczona i ciągła}, a $f:\R^d\to\R$
ograniczona i borelowska. Niech $\Delta=T/n$,
$t_k=k\Delta$. Wtedy
\[
\begin{split}
 (P_T^Vf)(x)&=\lim_{n\to\infty}
 \int_{(\R^d)^n}f(x_n)
 \exp\left(-\Delta\sum_{k=0}^{n-1}V(x_k)\right)\\
 &\hspace{2cm}\cdot\prod_{k=1}^n p_\Delta(x_{k-1},x_k)
 \,\dd\lambda_{dn}(x_1,\ldots,x_n),\qquad x_0=x.
\end{split}
\tag{15.10}\label{eq:path-time-slice}
\]
Dla ustalonych $x,y$ jądro \eqref{eq:path-FK-kernel}
jest granicą analogicznych całek po
$(\R^d)^{n-1}$ z $x_n=y$ (\emph{bez} końcowego
czynnika $f(x_n)$). W tych całkach iloczyn jąder
cieplnych pozostaje bez normalizacji; dopiero po
podzieleniu go przez $p_T(x,y)$ otrzymuje się gęstość
rozkładu punktów pośrednich mostu.
\end{twierdzenie}
\begin{proof}
Z Twierdzenia~\ref{th:path-cylinder} całka po
$(\R^d)^n$ jest równa
$\E_x[e^{-\Delta\sum_{k=0}^{n-1}V(X_{k\Delta})}f(X_T)]$.
Do zbieżności wykładnika nie potrzeba ciągłości $f$:
wartość $f(X_T)$ pozostaje taka sama dla wszystkich $n$.
Dla każdej ciągłej trajektorii $w$ jej obraz $w([0,T])$
jest zwarty, a $V$ jest na tym obrazie jednostajnie ciągła.
Zatem lewostronne sumy Riemanna zbiegają do
$\int_0^TV(w(s))\,\dd s$. Wykładniki są ograniczone
przez $e^{T\|V\|_\infty}$, więc twierdzenie o zbieżności
zdominowanej prowadzi do \eqref{eq:path-time-slice}.

Dla mostu \eqref{eq:bridge-density} daje rozkład
kolejnych położeń przed chwilą $T$; ten sam argument
Riemanna i zbieżności zdominowanej obowiązuje przy
każdym ustalonym $y$. Mnożąc przez $p_T(x,y)$,
otrzymujemy jądro \eqref{eq:path-FK-kernel}.
\end{proof}

\begin{uwaga}[Dlaczego wymagana jest ciągłość $V$?]
Dla ograniczonego potencjału borelowskiego
$\int_0^TV(X_s)\,\dd s$ nadal jest określona,
a twierdzenie o jądrze pozostaje prawdziwe.
Nie wolno jednak automatycznie twierdzić, że \emph{sumy
Riemanna przy ustalonej siatce} zbiegają dla każdej
trajektorii: funkcja $s\mapsto V(w(s))$ może być
nieciągła i niecałkowalna w sensie Riemanna.
Takie uogólnienie wymaga innej aproksymacji.
\end{uwaga}

\begin{center}
\begin{tikzpicture}[x=0.9cm,y=0.8cm,>=Latex]
\draw[->] (0,0)--(10.2,0) node[right] {$s$};
\draw[->] (0,-1.2)--(0,2.7) node[above] {$w(s)$};
\draw[thick,blue!65!black] plot[smooth] coordinates
 {(0,0.1) (0.6,0.5) (1.0,-0.05) (1.5,1.05) (2.0,0.55)
 (2.4,0.95) (3.1,0.3) (3.8,1.45) (4.2,0.85)
 (5.0,1.65) (5.8,0.95) (6.5,1.85) (7.0,0.8)
 (7.8,1.4) (8.3,0.75) (9.0,1.7) (9.8,1.15)};
\foreach \x in {0,2,4,6,8,10} {\draw[gray!70] (\x,-0.14)--(\x,0.14);}
\draw[dashed,gray!65] (2,-1.08)--(2,2.25);
\draw[dashed,gray!65] (4,-1.08)--(4,2.25);
\draw[dashed,gray!65] (6,-1.08)--(6,2.25);
\draw[dashed,gray!65] (8,-1.08)--(8,2.25);
\node[font=\small,align=center] at (5,-1.72)
 {Jedna trajektoria ciągła i kolejne chwile dyskretyzacji;\\
  całka po trajektoriach uśrednia funkcjonały po wszystkich takich drogach.};
\end{tikzpicture}
\end{center}

\section{Dwa zadania łączące probabilistykę z potencjałami}

\begin{zadanie}[Potencjał stały i ubytek masy]
Niech $V(x)\equiv v\ge0$. Wyznaczyć $P_t^V1$,
jądro i wyjaśnić sens probabilistyczny.

\emph{Rozwiązanie.} Ponieważ całka $\int_0^t V(X_s)\,\dd s=vt$
nie zależy od trajektorii, mamy
$(P_t^Vf)(x)=e^{-vt}\E_xf(X_t)$ i
$k_t^V(x,y)=e^{-vt}p_t(x,y)$. W szczególności
$(P_t^V1)(x)=e^{-vt}$. Ubytek masy interpretuje się
jako przeżycie cząstki zabijanej z intensywnością $v$;
nie jest to przejściowe jądro procesu konserwatywnego,
bo całka jądra po $y$ jest mniejsza niż $1$ dla $t>0$.
\end{zadanie}

\begin{zadanie}[Liczba odwiedzin a czas przebywania]
Niech $A\subset\R^d$ będzie borelowski,
$\lambda>0$ i $V=\lambda\ind{A}$. Wykazać, że
\[
 (P_T^V1)(x)=\E_x e^{-\lambda\,\ell_T(A)},
 \qquad \ell_T(A)=\int_0^T\ind{\{X_s\in A\}}\,\dd s.
\]

\emph{Rozwiązanie.} Podstawiamy $V$ w
\eqref{eq:path-FK-semigroup}. Zmienna $\ell_T(A)$
to \emph{czas okupacji} zbioru, a nie liczba jego odwiedzin.
Dla Brownowskiej trajektorii pojedyncze przecięcie granicy
zbioru może nie mieć sensownej liczności, lecz czas
okupacji pozostaje poprawnie określoną zmienną z przedziału
$[0,T]$. Z monotoniczności funkcji wykładniczej wynika
$0\le (P_T^V1)(x)\le1$, co odpowiada prawdopodobieństwu
przeżycia przy intensywności zabijania $\lambda$ wyłącznie
podczas pobytu w $A$.
\end{zadanie}

\section{Przejście do mechaniki kwantowej}

Wszystkie obiekty tego rozdziału są ściśle określonymi
całkami Lebesgue'a względem \emph{dodatnich} miar
probabilistycznych. W szczególności
$e^{-\int_0^T V(w(s))\,\dd s}$ jest dodatnią wagą dla
rzeczywistego $V$. W fizyce kwantowej pojawia się natomiast
czynnik $e^{iS[w]/\hbar}$, gdzie $S$ jest działaniem
klasycznym. W następnym rozdziale pokażemy, skąd bierze się
formalna analogia oraz dlaczego nie wolno zastępować
miary Wienera nieokreśloną „miarą fazową” bez oddzielnego
uzasadnienia matematycznego.


% ===========================================================================
% ROZDZIAŁ 16: całki Feynmana, Wick, propagator, oscylator
% ===========================================================================
